\documentclass[letterpaper,10pt]{texMemo}

\usepackage{graphicx}
\usepackage{listings}

\memoto{Dr. Mark Yoder}
\memofrom{Michael McDonald}
\memosubject{Problem Set 2 (8, 10, 10a, 11) Memo}
\memodate{\today}
\logo{\includegraphics[width=0.5\textwidth]{RoseLogo.png}}

\begin{document}
\maketitle

\section{Weekly Summary}

\noindent This week was devoted to GPIO, both using the Linux file system as
well as using C to access the files. Most of the exercises centered around the
idea of reading from files to gather input and then writing to files to produce
output. There was also a little in exercises 8 and 8a on the kernel and cross
compiling for ARMs. All source code is stored on my git repo. \\

\begin{center}
\begin{tabular}{|p{4cm}|l|l|p{8cm}|}
\hline
\textbf{Objective} & \textbf{Exercise} &\textbf{Status} & \textbf{Notes}\\
\hline
Compile the 3.8 Kernel & Exercise 8 & Completed & I didn't have to deal with
the ncurses failed dependency.\\
\hline
Install necessary cross compilation tools & Exercise 8a & Completed & \\
\hline
Read and Control GPIO's via the Linux File System & Exercise 10 & Completed &
Some I/O pins don't work well, possibly due to the presence of pull up/pull
down resistors and those being enabled.\\
\hline
Read analog input values from AIN channels & Exercise 10a & Completed & I
actually wish they output a straight numerical value, but I guess I can deal
with 1mV per division. It makes the math simpler, and is probably enough
resolution for most people.\\
\hline
Control GPIO's via C & Exercise 11 & Completed & This part was really cool. It
was nice to see the C code mesh with the real world (not just key inputs). \\
\hline
\end{tabular}
\end{center}

\subsection{Exercise Summary}

\subsubsection{Exercise 8}
\noindent I feel like I have ascended to a whole new plateau of nerddom after
having built the Linux Kernel. It is amazing. That's really all I have to say
about that. Everything went surprisingly smoothly. I didn't have to deal with
the ncurses library dependency issue, which was nice. Also, downloading
everything only took about 12 minutes, and compilation probably took about that
long as well. The video went a little fast to follow, I wish we could watch it
at about half speed.

\subsubsection{Exercise 8a}
\noindent Cross compilation went off without a hitch. Pretty cool to be able to
do that.

\subsubsection{Exercise 10}

\noindent Reading a switch can be accomplished by the following commands
\begin{verbatim}
 host$ sys/class/gpio: echo XX > export # add pin XX
 host$ sys/class/gpio: cd gpioXX # go into the gpio directory
 host$ sys/class/gpio/gpioXX: echo in > direction # make the pin an input
 host$ sys/class/gpio/gpioXX: cat value # read the value in from the port
\end{verbatim}

\begin{enumerate}
 \item What's the min and max voltage? Min: -100mV, Max 3.41V.
 \item What period is it? The period varies from 250ms to 260ms, with
occasional jumps as high as 350ms or 400ms.
 \item How close is it to 200ms? Off by about 50ms
 \item Why do they differ? Because the OS schedules it to run and occasionally
gives priority to other processes.
 \item Run htop and see how much processor you are using. 1-2\% of the CPU, and
less than 0.25\% of the RAM.
 \item Try different values for the sleep time (2nd argument). What's the
shortest
period you can get? Make a table of the values you try and the corresponding
period and processor usage. THe shortest period I can get is around 70 to 80ms.
 \item How stable is the period? The period is not very stable, it usually
varies by 10 to 20ms, and occasionally jumps 100 to 150ms higher.
 \item Try launching something like mplayer. How stable is the period? I
launched gedit, and it was very unstable when the app was launching, but it
stabilized close to what we saw earlier when it was just running in the
background. When it was in use, however, the period was much less stable, and
longer (300-400ms).
 \item Try cleaning up togglegpio.sh and removing unneeded lines. Does it impact
the period? Cleaning it up shaved off about 20-30ms, which is 10\% or so, so it
did a pretty good job in my opinion.
 \item togglegpio uses bash (first line in file). Try using sh. Is the period
shorter? Running it with sh seems to take about the same amount of time, though
it seems to be a little more stable.
 \item What's the shortest period you can get? sh could run much faster than
bash, and I saw periods as short 25-30ms.
\end{enumerate}

\begin{center}
\begin{tabular}{| l | l | l |}
\hline
\textbf{Time (s)} & \textbf{Period (ms)} & \textbf{CPU Load (\%)}\\
\hline
0.1 & 250-260 & 1-2\\
\hline
0.05 & 150-170 & 3-4\\
\hline
0.01 & 70-150  & 5-6\\ 
\hline
0.01 & 70-200 & 6-8\\
\hline
\end{tabular}
\end{center}

%\newpage
% Modify togglegpio.sh (call it toggleLED) to toggle the on-board LEDs. Can you
% get the LED to appear to dim by changing the duty cycle of the toggling?
%\lstinputlisting{toggleLED.sh}

%\newpage
% Display a count of how many times button has been pressed
%\lstinputlisting{countgpio.sh}

%\newpage
% Write a shell script that copies the value of gpio pin 7 (your switch) to gpio
% pin 60 (Your LED). How much CPU time does it take? What's the delay from the
% time the input changes until the output changes? How constant is the delay?
% 
% You may need to get a function generator and set it to a square wave (0 to 3.3V)
% and attach it to pin 7 and attached an oscilloscope to pin 60. 
%\lstinputlisting{controlLED.sh}

\noindent The controlLED shell script uses 6-8\% of the CPU, and less than
0.25\% of the memory, according to htop. The delay is about 20ms, and it remains
fairly constant.

\subsubsection{Exercise 10a}
\noindent The min and max value that I've been able to read are 0 to 1800,
and 12 bits is 4096. Therefore, it is reasonable to assume that AIN0-AIN6
output a 0.000 to 1.800 volt signal, where a difference by one is a 1mV
difference.

%\lstinputlisting{continuousADC.sh}

%\newpage
\subsubsection{Exercise 11}
% Compile and run it. Does the correct waveform appear on the oscilloscope?
% Use htop to measure the CPU usage.
% Try different periods. Make a chart of the input period vs. the measured
% period vs. CPU usage.
% How stable is the waveform? Look at several cycles. Are they always the same
% length. Explain. 

\noindent The correct waveform does appear on the scope. I set up a 50\% duty
cycle square wave with a period of 0.2 seconds, and I'm getting a 50\% duty
cycle square wave with a Frequency of 4.98-5.00 hz (or 200-201ms period), which
is fairly close to my 0.2 second desired period. The waveform is remarkably
stable, at least for relatively slow waveforms; however, for higher frequency
waveforms, they become less and less stable, and require more and more CPU
power to get the timing correct. Additionally, there are occasionally stretches
where the waveform is interrupted as the OS is servicing other requests. This
data is presented in the table below. \\

\begin{center}
\begin{tabular}{| l | l | l |}
\hline
\textbf{Desired Period (ms)} & \textbf{Actual Period (ms)} & \textbf{CPU Load
(\%)}\\
\hline
1000 & 1000 & 0 \\
\hline
100 & 100-101  & 0 \\
\hline
10 & 10.5-12.0 & 4 \\ 
\hline
1 & 1.25-2.0 & 14-16 \\
\hline
0.1 & 0.4-1 & 30-31 \\
\hline
\end{tabular}
\end{center}
% Modify togglegpio.c to take the gpio number.
% Presently one parameter specifies both the on and off time. Modify the code
% so the on and off times are controlled separately.
% Add in interrupt handler (see gpio-int-test.c) to trap ctrl-C and close
% things properly.
% What's the highest frequency you can see on the 'scope? 

%\newpage
%\lstinputlisting{gpio/togglegpio.c}
\noindent The highest frequency seen in my modified togglegpio was 2-3kHz,
however it was very unstable, and I believe the highest stable frequencies are
those less than 100Hz. \\

% Modify gpio-init-test.c to count the number of times the User button has
% been pressed. Set edge to only count releases.

%\newpage
%\lstinputlisting{gpio/gpio-int-test.c}
% Copy gpio-init-test.c to gpioThru.c and modify it to copy the value of one
% port to another. You'll have to add code to open a second gpio port for writing
% (check gpio_set_dir).
% gpio7 is attached to pin P9_42 on the Bone. Attach a function generator to
% this pin. Be sure to set the HiLevel to 3.3V and the LoLevel to 0. Does the
% output track the input?
% What's the highest frequency the output will track the input? What's the CPU
% usage?
% What's the delay from when the input changes to when the output changes? 

%\newpage
%\lstinputlisting{gpio/gpioThru.c}
\noindent The output tracks the input fairly well, usually about 0.2 to 0.3ms
behind, but sometimes stretches to 1-2ms. Depending on the input frequency, and
the other programs running on the bone, the signal gets less stable the higher
the frequency and the more programs running. It will track reliably up to 100
or 200 Hz, and will still track somewhat well to 1kHz. Above 1kHz however, it
really breaks down and by 2kHz it's almost totally unrecognizable. CPU usage is
nearly zero for frequencies below 10Hz, and they rise in the single digits
until about 250Hz, and by the time 1kHz rolls around, the program takes up
about 25\% of the CPU.

\subsection{GPIO input and Etch-a-Sketch}
\noindent The actual homework assignment was really easy to complete once the
exercises were completed. The etch-a-sketch took some time, but a majority of
that time was due to a bit of proper design, which paid off in the end, where I
literally combined my two final products (the etch-a-sketch and the four
input/output C program) and they just worked perfectly the first time.
Shocking, really. I finally feel like I'm getting a hang of physical computing,
since I feel like that was a non-trivial product that was created without
serious difficulty. I also really like being able to deal with physical input,
such as the switches that I wired up, and see that do something. Like the rest
of the exercises in this memo, all the source is available on my github repo. \\


\subsection{Takeaway Points}
\noindent The main points I took away from these exercises were the following:
\begin{itemize}
  \item Linux is impressively huge. I'm pretty amazed that people created it and
still work on it, and I hope to work on it as well.
  \item I think I have started to get a better feel for how computers have been
developed, interfacing all of the hardware I/O together to get say key input
and light output. I think it will be really cool to see the I2C LED screen
hooked up next week.
  \item I wish the exercises were renumbered to deal with dependencies (so that
exercise 8a came before exercise 8, for instance). It occasionally gets
confusing to be jumping around between exercises that are non-sequential to do
dependent tasks.
\end{itemize}

\end{document}
